\documentclass[11pt,a4paper]{article}
% preamble.tex — shared setup for human-verification documents
% Compile each document with:  tectonic <File>.tex

\usepackage[T1]{fontenc}
\usepackage[utf8]{inputenc}
\usepackage{lmodern}
\usepackage{amssymb,amsmath,amsthm}
\usepackage{mathtools}
\usepackage[dvipsnames]{xcolor}
\usepackage{listings}
\usepackage{hyperref}
\usepackage[margin=2.5cm]{geometry}
\usepackage{parskip}

% ── Colours ──────────────────────────────────────────────────────────
\definecolor{leanBg}{HTML}{F7F7F0}
\definecolor{haskellBg}{HTML}{F0F4F7}
\definecolor{keyword}{HTML}{7B3294}
\definecolor{comment}{HTML}{5A8C32}
\definecolor{noteFg}{HTML}{B35900}

% ── Lean 4 listings language ─────────────────────────────────────────
% NOTE: Lean listing blocks use ASCII approximations of unicode symbols
% (e.g. -> for arrow, o. for compose).  The literate key substitutes
% these with proper math-mode renderings in the PDF.
\lstdefinelanguage{Lean4}{
  morekeywords={def,theorem,structure,where,namespace,end,import,fun,
                let,have,show,by,exact,intro,cases,rfl,subst,obtain,
                funext,simp,rw,apply,variable,Type,Prop,instance,
                section,open,class,extends,deriving,sorry,noncomputable},
  sensitive=true,
  morecomment=[l]{--},
  morecomment=[n]{/-}{-/},
  morestring=[b]",
}

% ── Lean listing style ───────────────────────────────────────────────
\lstdefinestyle{lean}{
  language=Lean4,
  backgroundcolor=\color{leanBg},
  basicstyle=\small\ttfamily,
  keywordstyle=\color{keyword}\bfseries,
  commentstyle=\color{comment}\itshape,
  breaklines=true,
  frame=single,
  framerule=0.4pt,
  xleftmargin=4pt,
  xrightmargin=4pt,
  aboveskip=6pt,
  belowskip=4pt,
  literate=
    {->}{{$\to$}}2
    {=>}{{$\Rightarrow$}}2
    {<|}{{$\langle$}}1
    {|>}{{$\rangle$}}1
    {==>}{{$\Longrightarrow$}}3
    { o. }{{$\;\circ\;$}}3
    {forall}{{$\forall$}}1
    {:=}{{$\,:=\,$}}2
    {!=}{{$\neq$}}2
    {<=}{{$\leq$}}2
    {>=}{{$\geq$}}2,
}

% ── Haskell listing style ───────────────────────────────────────────
\lstdefinestyle{haskell}{
  language=Haskell,
  backgroundcolor=\color{haskellBg},
  basicstyle=\small\ttfamily,
  keywordstyle=\color{keyword}\bfseries,
  commentstyle=\color{comment}\itshape,
  breaklines=true,
  frame=single,
  framerule=0.4pt,
  xleftmargin=4pt,
  xrightmargin=4pt,
  aboveskip=6pt,
  belowskip=4pt,
  literate=
    {->}{{$\to$}}2
    {=>}{{$\Rightarrow$}}2,
}

% ── Theorem environments ────────────────────────────────────────────
\theoremstyle{definition}
\newtheorem{defn}{Definition}[section]
\newtheorem{thm}[defn]{Theorem}
\newtheorem{prop}[defn]{Proposition}

\theoremstyle{remark}
\newtheorem*{fidelitynote}{Fidelity note}

% ── Convenience macros ───────────────────────────────────────────────
\newcommand{\lean}[1]{\texttt{#1}}      % inline Lean identifier
\newcommand{\hask}[1]{\texttt{#1}}      % inline Haskell identifier
\newcommand{\Ext}{\mathrm{Ext}}
\newcommand{\Pos}{\mathrm{Pos}}
\newcommand{\Shape}{\mathrm{Shape}}


\title{Human Verification: \texttt{Containers.Thin}\\[4pt]
  \large Polytrees and Thinness of the Network Directed Container}
\author{Machine-generated from \texttt{lean/Thin.lean} for review by Neil Ghani}
\date{July 2026}

\begin{document}
\maketitle

\begin{abstract}
This document presents every definition and theorem \emph{statement} from the
Lean~4 module \lean{Containers.Thin} in three parallel forms: the original
Lean, a readable Haskell rendering, and precise mathematical English.  Proofs
are omitted --- they are machine-checked by Lean.  The purpose is to let a
human verifier confirm that the formalisation captures the intended
mathematics.
\end{abstract}

\tableofcontents
\bigskip

%===================================================================
\section{Overview}
%===================================================================

This file verifies Proposition~\emph{prop:tree} of the supply-chains paper: if
$G$ is a \textbf{polytree} (underlying undirected graph is a forest), then the
directed container $\mathrm{network}(G)$ is \textbf{thin} --- between any two
nodes there is at most one directed path.

The main results are:
\begin{itemize}
  \item \lean{thin\_of\_polytree}: in a polytree, any two directed paths
        between the same endpoints are equal (thinness = path uniqueness).
  \item \lean{endpoint\_equiv}: positions at $v$ in $\mathrm{network}(G)$
        biject with reachable nodes from $v$.
\end{itemize}

All definitions and proofs are in Lean~4 core; no Mathlib.

%===================================================================
\section{Reachability and path structure}
%===================================================================

% --- 1. Reachable -------------------------------------------------
\subsection{Definition: \lean{Digraph.Reachable}}

\begin{defn}[Reachability]
A vertex $w$ is \textbf{reachable} from $v$ in digraph $G$ if there exists a
directed path from $v$ to $w$:
\[
  \mathrm{Reachable}(G, v, w) \;:\!\iff\; \mathrm{Nonempty}(\mathrm{Path}(G, v, w)).
\]
\end{defn}

\paragraph{Lean.}
\begin{lstlisting}[style=lean]
def Reachable (G : Digraph) (v w : G.V) : Prop :=
  Nonempty (Path G v w)
\end{lstlisting}

\paragraph{Haskell.}
\begin{lstlisting}[style=haskell]
type Reachable g v w = Bool   -- True iff there is a path from v to w
reachable :: Digraph g => g -> Vertex -> Vertex -> Bool
reachable g v w = not (null (allPaths g v w))
\end{lstlisting}


% --- 2. Path.isTrivial --------------------------------------------
\subsection{Definition: \lean{Path.isTrivial}}

\begin{defn}[Trivial path]
A path $p$ is \textbf{trivial} (empty) when it was constructed by \lean{nil};
a \lean{cons} path is always non-trivial:
\[
  \mathrm{isTrivial}(\mathrm{nil}\,v) = \top, \qquad
  \mathrm{isTrivial}(\mathrm{cons}\,e\,p) = \bot.
\]
\end{defn}

\paragraph{Lean.}
\begin{lstlisting}[style=lean]
def Path.isTrivial : {v w : G.V} -> Path G v w -> Prop
  | _, _, .nil _ => True
  | _, _, .cons _ _ => False
\end{lstlisting}

\paragraph{Haskell.}
\begin{lstlisting}[style=haskell]
isTrivial :: Path -> Bool
isTrivial []      = True   -- nil path
isTrivial (_:_)   = False  -- cons path
\end{lstlisting}


% --- 3. Acyclic ---------------------------------------------------
\subsection{Definition: \lean{Digraph.Acyclic}}

\begin{defn}[Acyclic graph]
A digraph $G$ is \textbf{acyclic} if every path from a vertex to itself is
trivial (empty):
\[
  \mathrm{Acyclic}(G) \;:\!\iff\; \forall v : V,\;\forall p : \mathrm{Path}(G, v, v),\;
  p.\mathrm{isTrivial}.
\]
\end{defn}

\paragraph{Lean.}
\begin{lstlisting}[style=lean]
def Acyclic (G : Digraph) : Prop :=
  forall (v : G.V) (p : Path G v v), p.isTrivial
\end{lstlisting}

\paragraph{Haskell.}
\begin{lstlisting}[style=haskell]
isAcyclic :: Digraph g => g -> Bool
isAcyclic g = all (\v -> isTrivial (shortestCycle g v)) (vertices g)
-- Equivalently: the graph has no directed cycles.
\end{lstlisting}


% --- 4. EdgeDisjoint ----------------------------------------------
\subsection{Definition: \lean{Digraph.EdgeDisjoint}}

\begin{defn}[Edge-disjoint reachability]
A digraph $G$ has \textbf{edge-disjoint reachability} if for any two distinct
edges $e_1 \neq e_2$ with the same source, no vertex is reachable from both
targets:
\[
  \mathrm{EdgeDisjoint}(G) \;:\!\iff\; \forall e_1, e_2 : E,\;
  \mathrm{src}(e_1) = \mathrm{src}(e_2) \Rightarrow e_1 \neq e_2 \Rightarrow
  \forall w,\; \neg(\mathrm{Reachable}(G, \mathrm{tgt}(e_1), w) \wedge
  \mathrm{Reachable}(G, \mathrm{tgt}(e_2), w)).
\]
\end{defn}

\paragraph{Lean.}
\begin{lstlisting}[style=lean]
def EdgeDisjoint (G : Digraph) : Prop :=
  forall (e1 e2 : G.E), G.src e1 = G.src e2 -> e1 != e2 ->
    forall w, -(Reachable G (G.tgt e1) w /\ Reachable G (G.tgt e2) w)
\end{lstlisting}

\paragraph{Haskell.}
\begin{lstlisting}[style=haskell]
edgeDisjoint :: Digraph g => g -> Bool
edgeDisjoint g = and
  [ null (reachableSet g (tgt e1) `intersect` reachableSet g (tgt e2))
  | e1 <- edges g, e2 <- edges g, e1 /= e2, src e1 == src e2 ]
\end{lstlisting}


% --- 5. Polytree --------------------------------------------------
\subsection{Structure: \lean{Digraph.Polytree}}

\begin{defn}[Polytree]
A \textbf{polytree} is a digraph that is both acyclic and edge-disjoint:
\[
  \mathrm{Polytree}(G) \;:\!\iff\; \mathrm{Acyclic}(G) \wedge \mathrm{EdgeDisjoint}(G).
\]
\end{defn}

\paragraph{Lean.}
\begin{lstlisting}[style=lean]
structure Polytree (G : Digraph) : Prop where
  acyclic : Acyclic G
  edgeDisjoint : EdgeDisjoint G
\end{lstlisting}

\paragraph{Haskell.}
\begin{lstlisting}[style=haskell]
data Polytree g = Polytree
  { acyclic      :: Bool   -- isAcyclic g
  , edgeDisjoint :: Bool   -- edgeDisjoint g
  }
\end{lstlisting}


%===================================================================
\section{Auxiliary lemmas and path decomposition}
%===================================================================

% --- 6. acyclic_no_cons -------------------------------------------
\subsection{Lemma: \lean{acyclic\_no\_cons}}

\begin{prop}[Acyclicity excludes cons self-loops]
In an acyclic graph, if $e : E$ is an edge, there is no directed path from
$\mathrm{tgt}(e)$ back to $\mathrm{src}(e)$:
\[
  \mathrm{Acyclic}(G) \Rightarrow \neg\,\mathrm{Path}(G,\, \mathrm{tgt}(e),\, \mathrm{src}(e)).
\]
\end{prop}

\paragraph{Lean.}
\begin{lstlisting}[style=lean]
theorem acyclic_no_cons (hac : Acyclic G) (e : G.E) (p : Path G (G.tgt e) (G.src e)) :
    False
\end{lstlisting}

\paragraph{Haskell.}
\begin{lstlisting}[style=haskell]
-- In an acyclic graph, no back-edge path exists:
-- acyclic_no_cons :: Acyclic g -> Edge -> Path (tgt e) (src e) -> Void
\end{lstlisting}


% --- 7. PathView --------------------------------------------------
\subsection{Inductive type: \lean{PathView}}

\begin{defn}[Path view]
A \textbf{path view} from $v$ to $w$ is either:
\begin{itemize}
  \item a \textbf{trivial} view with a proof $h : v = w$, or
  \item a \textbf{step}: an edge $e$ with $\mathrm{src}(e) = v$, followed by a
        path $\mathrm{tail} : \mathrm{Path}(G, \mathrm{tgt}(e), w)$.
\end{itemize}
This packages the index equalities explicitly to enable clean pattern matching
without triggering dependent-elimination issues.
\end{defn}

\paragraph{Lean.}
\begin{lstlisting}[style=lean]
inductive PathView (G : Digraph) (v w : G.V) : Type
  | trivial (h : v = w) : PathView G v w
  | step (e : G.E) (he : G.src e = v) (tail : Path G (G.tgt e) w)
      : PathView G v w
\end{lstlisting}

\paragraph{Haskell.}
\begin{lstlisting}[style=haskell]
data PathView v w
  = Trivial    -- v = w
  | Step Edge (Path w)   -- an edge followed by a path to w
\end{lstlisting}


% --- 8. Path.view / PathView.toPath / Path.view_toPath -----------
\subsection{Definitions and theorem: \lean{Path.view}, \lean{PathView.toPath}, \lean{Path.view\_toPath}}

\begin{defn}[View of a path]
Every path can be viewed as trivial or a step:
\begin{align*}
  \mathrm{view}(\mathrm{nil}\,v) &= \mathrm{trivial}(\mathrm{rfl}), \\
  \mathrm{view}(\mathrm{cons}\,e\,p) &= \mathrm{step}(e, \mathrm{rfl}, p).
\end{align*}
\end{defn}

\begin{defn}[Path from a view]
Recover the original path from its view:
\begin{align*}
  \mathrm{toPath}(\mathrm{trivial}\,h) &= h \triangleright \mathrm{nil}, \\
  \mathrm{toPath}(\mathrm{step}\,e\,h_e\,p) &= h_e \triangleright \mathrm{cons}(e, p).
\end{align*}
\end{defn}

\begin{thm}[View round-trips]
For any path $p$, $\mathrm{toPath}(\mathrm{view}(p)) = p$.
\end{thm}

\paragraph{Lean.}
\begin{lstlisting}[style=lean]
def Path.view : {v w : G.V} -> Path G v w -> PathView G v w
  | _, _, .nil v => .trivial rfl
  | _, _, .cons e p => .step e rfl p

def PathView.toPath : {v w : G.V} -> PathView G v w -> Path G v w
  | _, _, .trivial h => h [> .nil _
  | _, _, .step e he p => he [> .cons e p

theorem Path.view_toPath : {v w : G.V} -> (p : Path G v w) ->
    p.view.toPath = p
\end{lstlisting}

\paragraph{Haskell.}
\begin{lstlisting}[style=haskell]
viewPath :: Path v w -> PathView v w
viewPath []     = Trivial
viewPath (e:es) = Step e es

toPath :: PathView v w -> Path v w
toPath Trivial       = []
toPath (Step e es) = e : es

-- prop_view_roundtrip p = toPath (viewPath p) == p
\end{lstlisting}


%===================================================================
\section{Main results}
%===================================================================

% --- 9. path_unique_of_polytree -----------------------------------
\subsection{Theorem: \lean{path\_unique\_of\_polytree}}

\begin{thm}[Path uniqueness in polytrees]
In a polytree $G$, any two directed paths between the same endpoints are equal:
\[
  \mathrm{Polytree}(G) \;\Rightarrow\;
  \forall v, w : V,\; \forall p, q : \mathrm{Path}(G, v, w),\; p = q.
\]
\emph{Proof strategy}: structural induction on $p$, decomposing $q$ via
\lean{PathView} to avoid dependent-elimination issues.
\begin{itemize}
  \item \lean{nil}/\lean{trivial}: both paths are \lean{nil}.
  \item \lean{nil}/\lean{step}: the step creates a directed cycle, contradicting
        acyclicity.
  \item \lean{cons}/\lean{trivial}: symmetric cycle contradiction.
  \item \lean{cons} $e_1$ / \lean{step} $e_2$: if $e_1 = e_2$, apply the
        induction hypothesis to the tails; if $e_1 \neq e_2$, edge-disjointness
        is violated.
\end{itemize}
\end{thm}

\paragraph{Lean.}
\begin{lstlisting}[style=lean]
theorem path_unique_of_polytree (ht : Polytree G) :
    {v w : G.V} -> (p q : Path G v w) -> p = q
\end{lstlisting}

\paragraph{Haskell.}
\begin{lstlisting}[style=haskell]
-- prop_unique :: Polytree g -> Path v w -> Path v w -> Bool
-- prop_unique ht p q = p == q
--
-- In a polytree, all paths between the same endpoints are identical.
\end{lstlisting}


% --- 10. thin_of_polytree -----------------------------------------
\subsection{Theorem: \lean{thin\_of\_polytree}}

\begin{thm}[Thinness of the network container]
If $G$ is a polytree, then for any two vertices $v$ and $w$, the type
$\mathrm{Path}(G, v, w)$ is a subsingleton (has at most one element):
\[
  \mathrm{Polytree}(G) \;\Rightarrow\;
  \forall v, w : V,\; \mathrm{Subsingleton}(\mathrm{Path}(G, v, w)).
\]
This is the thinness condition on the directed container $\mathrm{network}(G)$:
between any two nodes there is at most one path.
\end{thm}

\paragraph{Lean.}
\begin{lstlisting}[style=lean]
theorem thin_of_polytree (ht : Polytree G) (v w : G.V) :
    Subsingleton (Path G v w)
\end{lstlisting}

\paragraph{Haskell.}
\begin{lstlisting}[style=haskell]
-- prop_thin :: Polytree g -> Vertex -> Vertex -> Bool
-- prop_thin ht v w = length (allPaths g v w) <= 1
\end{lstlisting}


%===================================================================
\section{Position bijection}
%===================================================================

% --- 11. Equiv' ---------------------------------------------------
\subsection{Structure: \lean{Equiv'}}

\begin{defn}[Equivalence (Mathlib-free)]
An \textbf{equivalence} between types $\alpha$ and $\beta$ is a pair of inverse
maps:
\[
  \mathrm{Equiv}'(\alpha, \beta) \;=\;
  \{ (f, g, l, r) \mid f : \alpha \to \beta,\; g : \beta \to \alpha,\;
    l : g \circ f = \mathrm{id},\; r : f \circ g = \mathrm{id} \}.
\]
\end{defn}

\paragraph{Lean.}
\begin{lstlisting}[style=lean]
structure Equiv' (alpha beta : Type) where
  toFun    : alpha -> beta
  invFun   : beta -> alpha
  left_inv  : forall a, invFun (toFun a) = a
  right_inv : forall b, toFun (invFun b) = b
\end{lstlisting}

\paragraph{Haskell.}
\begin{lstlisting}[style=haskell]
data Equiv a b = Equiv
  { toFun    :: a -> b
  , invFun   :: b -> a
  -- left_inv:  invFun . toFun = id
  -- right_inv: toFun . invFun = id
  }
\end{lstlisting}


% --- 12. Reachable.path -------------------------------------------
\subsection{Definition: \lean{Reachable.path}}

\begin{defn}[Choosing a path witness]
Given a reachability proof $h_r : \mathrm{Reachable}(G, v, w)$, choose a
path by classical choice:
\[
  h_r.\mathrm{path} \;=\; \mathrm{Classical.choice}(h_r).
\]
This is \lean{noncomputable} (no constructive witness is extracted).
\end{defn}

\paragraph{Lean.}
\begin{lstlisting}[style=lean]
noncomputable def Reachable.path {v w : G.V}
    (hr : Reachable G v w) : Path G v w :=
  Classical.choice hr
\end{lstlisting}

\paragraph{Haskell.}
\begin{lstlisting}[style=haskell]
-- Non-constructive: pick any path given a reachability certificate.
-- In Haskell we would just compute one (e.g. BFS).
reachablePath :: Reachable -> Path
reachablePath = head . allPaths  -- pick arbitrarily
\end{lstlisting}

\begin{fidelitynote}
In Lean, \lean{noncomputable} signals that the definition uses
\lean{Classical.choice} and cannot be extracted to code.  In Haskell (or an
actual implementation) a constructive path would be computed by BFS or DFS.
\end{fidelitynote}


% --- 13. endpoint_equiv -------------------------------------------
\subsection{Definition: \lean{endpoint\_equiv}}

\begin{defn}[Positions biject with reachable nodes]
In a polytree $G$, positions at $v$ in $\mathrm{network}(G)$ --- which are
pairs $(w, p)$ with $p : \mathrm{Path}(G, v, w)$ --- biject with reachable
nodes from $v$:
\[
  \mathrm{endpoint\_equiv}(\mathrm{ht}, v) \;:\;
  \Bigl(\Sigma_{w : V}\, \mathrm{Path}(G, v, w)\Bigr)
  \;\simeq\;
  \{w : V \mid \mathrm{Reachable}(G, v, w)\}.
\]
\begin{itemize}
  \item Forward: $(w, p) \mapsto (w, \langle p \rangle)$ (pack into a
        reachability proof).
  \item Backward: $(w, h_r) \mapsto (w, h_r.\mathrm{path})$ (choose the
        unique path via \lean{path\_unique\_of\_polytree}).
\end{itemize}
This is \lean{noncomputable} (it uses \lean{Classical.choice} via
\lean{Reachable.path}).
\end{defn}

\paragraph{Lean.}
\begin{lstlisting}[style=lean]
noncomputable def endpoint_equiv (ht : Polytree G) (v : G.V) :
    Equiv' ((w : G.V) * Path G v w) {w : G.V // Reachable G v w} where
  toFun p := <|p.1, <|p.2|>|>
  invFun w := <|w.1, w.2.path|>
  left_inv p := ...
  right_inv w := ...
\end{lstlisting}

\paragraph{Haskell.}
\begin{lstlisting}[style=haskell]
endpoint_equiv :: Polytree g -> Vertex
               -> Equiv (Vertex, Path) (Vertex, Bool {- reachable -})
endpoint_equiv ht v = Equiv
  { toFun  = \(w, p) -> (w, True)
  , invFun = \(w, _) -> (w, uniquePath v w)  -- unique by polytree assumption
  }
\end{lstlisting}

\begin{fidelitynote}
The Lean statement uses a dependent $\Sigma$-type for positions and a
subtype for reachable nodes.  Haskell approximates these with pairs and a
boolean flag.  The key mathematical content --- the bijection between paths
and reachable nodes in a polytree --- is preserved.
\end{fidelitynote}


%===================================================================
\section{Translation conventions}
%===================================================================

\begin{enumerate}
\item \textbf{Classical choice.}  \lean{endpoint\_equiv} uses
  \lean{Classical.choice}; it is \lean{noncomputable}.  In a constructive
  or implementation setting, one would compute a path by graph search.

\item \textbf{Dependent types for positions.}  Positions in
  $\mathrm{network}(G)$ are $\Sigma_{w}\,\mathrm{Path}(G, v, w)$; the
  endpoint $w$ is \emph{determined} by the path.  Haskell approximates this
  with a plain pair.

\item \textbf{PathView.}  The \lean{PathView} type exists solely to avoid
  dependent-elimination failures when matching on the second path in the
  uniqueness proof.  It has no mathematical content beyond providing a clean
  view of the path type.

\item \textbf{Proof bodies omitted.}  Only statements are translated.
\end{enumerate}

\end{document}
